% Generated draft: requires mathematical and source review.
\documentclass[11pt]{article}
\usepackage[margin=25mm]{geometry}
\usepackage{fontspec}
\setmainfont{Latin Modern Roman}
\usepackage{amsmath,amssymb,mathtools,unicode-math}
\setmathfont{Latin Modern Math}
\usepackage{enumitem,needspace,xurl,hyperref}
\setlength{\parindent}{0pt}
\setlength{\parskip}{5pt}
\setlength{\emergencystretch}{3em}
\newcommand{\N}{\mathbb{N}}
\newcommand{\Z}{\mathbb{Z}}
\newcommand{\Q}{\mathbb{Q}}
\newcommand{\R}{\mathbb{R}}
\newcommand{\C}{\mathbb{C}}
\newcommand{\F}{\mathbb{F}}
\newcommand{\A}{\mathbb{A}}
\newcommand{\PP}{\mathbb{P}}
\newcommand{\E}{\mathbb{E}}
\newcommand{\eps}{\varepsilon}
\newcommand{\dd}{\,\mathrm{d}}
\begin{document}
\section*{UnsolvedMath ID: 800828}
\textbf{Source problem code:} catalogue entry
\subsection{Definitions and mathematical statement}
Let \(M\) be a compact manifold of dimension \(2n\), and let \(F_1,\ldots,F_n\) be the first integrals of the system described in the source. Write \(F=(F_1,\ldots,F_n)\). To interpret the assumptions precisely, one must specify the dynamical evolution, the bracket used for involution, the differentiability class, and the meaning of almost everywhere. If \(\phi_t\) denotes the evolution, \(\{\cdot,\cdot\}\) its specified bracket, and \(\mu\) the measure intended by that phrase, the supplied conditions take the form
\[
F_i\circ\phi_t=F_i\quad\text{wherever the evolution is defined},\qquad
\{F_i,F_j\}=0\quad(1\leq i,j\leq n),
\]
\[
\operatorname{rank}(\dd F_x)=n\qquad\text{for \(\mu\)-almost every }x\in M.
\]
Here these symbols express the source assumptions conditionally; the source does not specify a symplectic or Poisson structure, a measure, or a regularity class.

Define the singular locus of the integral map by
\[
\Sigma=\{x\in M:\operatorname{rank}(\dd F_x)<n\}.
\]
The set \(\Sigma\) need not itself be a submanifold. To speak of dynamics restricted to a submanifold \(S\subseteq\Sigma\), that submanifold must be invariant under the evolution; the source does not state this requirement explicitly.

\textbf{Research question (Khesin, Misiolek, Modin, Problem 27).} What restrictions does the stated integrability impose on the dynamics along singular submanifolds where the integrals become dependent? How complicated can such dynamics be, and can chaotic dynamics occur on some tiny invariant submanifold contained in \(\Sigma\)? The source does not define either chaotic dynamics or tiny, and consequently does not specify a unique existence assertion or classification target. A precise theorem formulation requires these meanings and the underlying geometric and regularity assumptions to be fixed.
\subsection{Short English statement}
The problem asks how complicated the dynamics of an integrable system on a compact \(2n\)-dimensional manifold can become where its \(n\) commuting first integrals lose independence. It asks in particular whether chaos can occur on a tiny singular submanifold. The geometric setting, regularity, meaning of almost everywhere, and meanings of chaos and tiny remain unspecified.
\subsection{Sources}
\begin{itemize}
\item Supplied local UnsolvedMath registry record: permanent ID \(800828\), exact source problem code \texttt{catalogue entry}; collection: AMR Open Problem Lists. The supplied registry labels the problem open; this label is not an independently verified mathematical or literature assessment.
\item Khesin, Misiolek, Modin, \emph{Geometric Hydrodynamics in Open Problems} (2022), Problem 27. Supplied source URL: \url{https://arxiv.org/abs/2205.01143}. The record attributes the problem to these authors and gives proposed year 2022.
\item Supplied provenance reports access on 2026-10-16 and extraction by \texttt{arxiv\_latex}. This is supplied provenance, not a claim of independent source access or verification. The registry record was created and updated at \texttt{2026-10-16T14:00:00Z}; its supplied upstream revision is \texttt{d371ee917ac9647fbb06cd6fc529deec6f3dda82}.
\end{itemize}
\end{document}
